% Generated by analysis_provenance/make_supp_sections.py from the saved outcomes. Do not edit.
\section{Second sea-clutter radar: NetRAD (R15)}\label{sup:r15}
Frozen protocol \texttt{r15/PROTOCOL\_R15.md} (sha256 945c5e51\ldots). It was written before any statistic of the NetRAD samples was computed. Before that point only the following had been read: file sizes, array shapes, and the providers' scripts and log. In addition, two of the providers' own quick-look images were viewed while the dataset was being located. Expectations and verdicts are in \texttt{r15/RESULTS\_R15.md} (12 of 15 met), and deviations in \texttt{r15/DEVIATIONS\_R15.md}. The deviations are a corrected site description (the Atlantic side of the Cape Peninsula, not False Bay) and run mechanics; none changes a method, setting or expectation.

\emph{Data.} Node 3 (monostatic) of the NetRAD system: seven HH and seven VV recordings of 130{,}000 pulses at 1~kHz. The clutter cells are the providers' per-recording range windows (\texttt{NR\_Range\_Selection\_June9.m}). Each cell has its complex mean removed and is scaled to unit power. The cross-polarized recordings were set aside before any data were read. All methods are the frozen R12 and R14 code. The only change is that the ANMF secondary cells are restricted to those two to four cells away, which gives 15--30 vectors (20--35 on IPIX). The AR(1) clamp $|r|\le0.98$ binds in 12 of the 14 recordings.

\emph{Unmet expectations of R15.}
\begin{itemize}
\item \textbf{Calibration at $10^{-4}$ (N1).} At least 12 of 13 conformal ratios were expected within $[0.5,2]$ at $10^{-4}$; 7 were. Outside: Clip-OS1 (0.39), D-IN4 (4.26), NA4 (2.08), P-ANMF (2.33), P-ANMF(bin) (2.52), PAMF-H(w) (2.18). At $10^{-3}$ all 13 are within $[0.5,2]$, and the 12 non-saturating statistics lie within 1.01--1.36 times design.
\item \textbf{Pivotality ranking (N4).} IN-ARCP's quintile spread at $10^{-2}$ was expected to be at most 2 and below CA16's. It is 2.40 against 1.51. In the same analysis on IPIX it is 1.31 against 1.06, so this expectation had no support on IPIX either.
\item \textbf{Flagged pulses (W2).} Non-coherent integration with a clean threshold was expected to exceed $1.5\alpha$ in segments with flagged pulses. It stays at 1.08 times design.
\end{itemize}
Expectation N5 (PAMF-H within 0.5~dB of the best dwell detector) was met, but it is uninformative: every whitened dwell detector reaches $P_{\rm d}=0.5$ at the lowest SCR tested.

\emph{Exploratory (post hoc, after the outcomes and the claim audit).} Two scripts look at where the conformal exceedances at $10^{-4}$ fall: \texttt{r15/explore\_r15\_hotspots.py} and \texttt{r15/explore\_r15\_attribution.py}.
\begin{itemize}
\item \textbf{Integration and PAMF-H.} Most of their excess comes from two VV units (1450 and 1501, assignment 3). Their exceedances fall within one to three seconds of the test third and coincide with flagged pulses, while their calibration thirds contain no flagged pulse. Without these two units the pooled ratios are 1.56 (integration) and 1.33 (PAMF-H).
\item \textbf{P-ANMF statistics.} Their excess comes from one HH unit (1128, assignment 3) with frequent flagged pulses in both thirds. Fewer than half of its exceedances coincide with flagged pulses.
\item \textbf{Unexplained.} The flagged pulses do not explain the excess of the noise-aware score (2.08). Nor do they explain an excess of integration of more than five times design in one unit that has no flagged pulse at all (1508, assignment 2).
\end{itemize}
\textbf{What the flags are is not established.} \texttt{r15/explore\_r15\_flag\_extent.py} counts, in the raw recordings, the range bins above ten times their median:
\begin{itemize}
\item a flagged HH pulse has a median of about 108 such bins, of 1024, and 85--91\% of them lie outside the clutter window;
\item unflagged HH pulses with an elevated window bin have about 80.
\end{itemize}
The flags are therefore the upper tail of a phenomenon that extends over range. They may be wideband interference: the providers remove WiFi pulses in their own processing of the bistatic nodes. They may equally be range-extended HH sea spikes. The main text therefore calls them flagged pulses, not interference.

In the clutter cells, uncertified non-coherent integration runs at 1.05 times design at $10^{-2}$ and 1.20 at $10^{-3}$ over all test segments of the 7 HH recordings, the only ones with at least 0.5\% flagged pulses (VV: at most 0.1\%). Segments without a flagged pulse give 1.04 and 1.17; segments with flagged pulses give at most 1.08 and 1.33, with wide intervals at $10^{-3}$ (Part W below). A certificate whose hit model must cover these frequent, bursty masks is nearly vacuous ($P_{\rm d}=0.04$). Blanking with mask-matched calibration keeps $P_{\rm d}=0.92$ at 0.57 times design.

\noindent Output of \texttt{r15/analyze\_r15.py} (Parts L, A, G, I, K/X/B and W; intervals resample recordings):
\VerbatimInput[fontsize=\tiny]{data/r15_summary.txt}
\noindent Exploratory hotspot analysis (\texttt{r15/explore\_r15\_hotspots.py}):
\VerbatimInput[fontsize=\tiny]{data/r15_hotspots.txt}
\noindent Exploratory attribution (\texttt{r15/explore\_r15\_attribution.py}): flagged-pulse fraction in each unit's calibration and test thirds, the per-unit ratios at $10^{-4}$, and the pooled ratios without the units named above:
\VerbatimInput[fontsize=\tiny]{data/r15_attribution.txt}
\noindent Exploratory range extent of the flagged pulses (\texttt{r15/explore\_r15\_flag\_extent.py}):
\VerbatimInput[fontsize=\tiny]{data/r15_flag_extent.txt}
